\documentclass[a4paper,12pt]{article}

\usepackage[italian]{babel}
\usepackage[utf8]{inputenc}
\usepackage{geometry}

\geometry{margin=2.5cm}

\title{CEREALEHUB \\ Sistema di Gestione Magazzino}
\author{Progetto sviluppato in Java}
\date{\today}

\begin{document}

\maketitle

\section{Introduzione}
Il progetto \textbf{CerealeHub} ha lo scopo di sviluppare un sistema software per la gestione di un magazzino logistico. 
Il sistema consente di gestire pacchi, utenti e corrieri, garantendo un'organizzazione efficiente delle spedizioni.

L'applicazione è sviluppata in \textbf{Java}, utilizzando i principi della programmazione orientata agli oggetti (OOP), 
senza l'utilizzo di API esterne.

\section{Obiettivi}
Gli obiettivi principali del sistema sono:
\begin{itemize}
    \item Gestione strutturata dei pacchi
    \item Differenziazione tra utenti base e premium
    \item Assegnazione automatica dei pacchi ai corrieri
    \item Simulazione del processo di spedizione
\end{itemize}

\section{Modello dei dati}

\subsection{Classe Pacco}
La classe \texttt{Pacco} rappresenta un'unità di spedizione.

Attributi:
\begin{itemize}
    \item \texttt{id} (int)
    \item \texttt{peso} (double)
    \item \texttt{destinazione} (String)
    \item \texttt{tipo} (enum: BASE, PREMIUM)
    \item \texttt{stato} (enum: IN\_MAGAZZINO, IN\_CONSEGNA, CONSEGNATO)
\end{itemize}

Metodi principali:
\begin{itemize}
    \item \texttt{cambiaStato()}
    \item \texttt{toString()}
\end{itemize}

\subsection{Classe Utente}
La classe \texttt{Utente} rappresenta i clienti.

Attributi:
\begin{itemize}
    \item \texttt{id}
    \item \texttt{nome}
    \item \texttt{tipo} (BASE o PREMIUM)
    \item \texttt{listaPacchi} (ArrayList di Pacco)
\end{itemize}

Metodi:
\begin{itemize}
    \item \texttt{aggiungiPacco(Pacco p)}
    \item \texttt{visualizzaPacchi()}
\end{itemize}

\subsection{Classe Corriere}
La classe \texttt{Corriere} rappresenta chi si occupa della consegna.

Attributi:
\begin{itemize}
    \item \texttt{id}
    \item \texttt{nome}
    \item \texttt{zona}
    \item \texttt{pacchiAssegnati} (ArrayList)
\end{itemize}

Metodi:
\begin{itemize}
    \item \texttt{assegnaPacco(Pacco p)}
    \item \texttt{consegnaPacchi()}
\end{itemize}

\subsection{Classe Magazzino}
La classe \texttt{Magazzino} gestisce l'intero sistema.

Attributi:
\begin{itemize}
    \item Lista pacchi
    \item Lista utenti
    \item Lista corrieri
\end{itemize}

Metodi:
\begin{itemize}
    \item \texttt{aggiungiPacco()}
    \item \texttt{registraUtente()}
    \item \texttt{aggiungiCorriere()}
    \item \texttt{assegnaPacchi()}
\end{itemize}

\section{Logica del sistema}

Il sistema implementa una logica leggermente più avanzata:

\begin{itemize}
    \item I pacchi \textbf{premium} vengono assegnati con priorità ai corrieri
    \item Ogni corriere ha un limite massimo di pacchi trasportabili
    \item I pacchi cambiano stato durante il ciclo di vita:
    \begin{itemize}
        \item In magazzino
        \item In consegna
        \item Consegnato
    \end{itemize}
\end{itemize}

\section{Strutture dati utilizzate}
Nel progetto vengono utilizzate:
\begin{itemize}
    \item \textbf{ArrayList} per la gestione dinamica degli oggetti
    \item \textbf{Enum} per rappresentare tipi e stati
    \item \textbf{Classi e oggetti} per modellare il dominio
\end{itemize}

\section{Esempio di flusso operativo}
\begin{enumerate}
    \item Un utente viene registrato
    \item L'utente inserisce un pacco
    \item Il pacco viene salvato nel magazzino
    \item Il sistema assegna automaticamente il pacco a un corriere
    \item Il corriere consegna il pacco aggiornandone lo stato
\end{enumerate}

\section{Possibili miglioramenti}
\begin{itemize}
    \item Introduzione di file per la persistenza dei dati
    \item Interfaccia grafica (GUI)
    \item Sistema di priorità più complesso
    \item Statistiche sulle consegne
\end{itemize}

\section{Conclusione}
Il sistema CerealeHub rappresenta una simulazione realistica di un sistema logistico. 
L'utilizzo di Java e dei principi OOP permette una struttura modulare e facilmente estendibile, 
rendendo il progetto adatto a evoluzioni future.

\end{document}